% ЕМТ 2023, XII клас — точен LaTeX препис от официалните файлове в Klasirane.com.
% Условия: https://klasirane.com/api/competitions/EMT/All/2023/12%20%D0%BA%D0%BB/probs
% SHA-256 (условия): 583f840deeeef77234080964f64ecb7ae6cbd5ae28898478f3ed8b306a185a84
% Решения: https://klasirane.com/api/competitions/EMT/All/2023/12%20%D0%BA%D0%BB/sol
% SHA-256 (решения): 4822dbfc17fbe39e12cc11d8d9d97d3e738f6d43b72b6f2b58a194bad5b7879a
% Точковите оценки, правописът, формулировките, липсващите символи и видимите печатни особености са запазени от източника.
% В официалните файлове няма отделна чертежна фигура; не е добавяна такава.

Задача 12.1. Редицата $\{x_n\}_{n=0}^{\infty}$ е зададена чрез равенствата:
$$
\begin{aligned}
x_0&=1,\\
x_{n+1}&=\sin x_n+\frac{\pi}{2}-1\text{ за }n\geq0.
\end{aligned}
$$
Да се докаже, че редицата $\{x_n\}_{n=0}^{\infty}$ е сходяща и да се определи границата ѝ.

Решение. Първо ще докажем, че редицата е строго растяща. Да отбележим, че за всяко $n$ имаме $x_n\leq\frac{\pi}{2}$, т.к $\sin x\leq1$ за всяко $x\in\mathbb R$. Освен това имаме, че функцията $f(x)=\sin x-x$ е намаляваща за $x\in\mathbb R$, т.к $f'(x)=\cos x-1\leq0$ за всяко $x\in\mathbb R$.

Следователно за $x\in\left(-\infty,\frac{\pi}{2}\right)$ имаме $f(x)\geq\sin\frac{\pi}{2}-\frac{\pi}{2}$. Така получаваме, че $\sin x_n\geq x_n+1-\frac{\pi}{2}$, което е еквивалентно на $x_{n+1}\geq x_n$. Следователно редицата $(x_n)_{n=1}^{\infty}$ е растяща и т.к тя е ограничена следва, че е сходяща. Ако $l$ е нейната граница, то за $l$ е изпълнено, че $l=\sin l+\frac{\pi}{2}-1$, т.е $f(l)=f\left(\frac{\pi}{2}\right)$. Функцията $f$ е намаляваща, което означава, че $l=\frac{\pi}{2}$.

Оценяване. (6 точки) 1т. за доказателство, че $x_n\leq\pi/2$ за всяко $n$, 3т. за доказателство, че $x_n$ е растяща и 2т. за довършване.

Задача 12.2. Даден е остроъгълен и разностранен триъгълник $ABC$. Вписаната в триъгълник $ABC$ окръжност с център $I$ допира страните $BC,AC$ и $AB$ съответно в точките $D,E$ и $F$. Окръжността с център $C$ и радиус $CE$ пресича за втори път правата $EF$ в точка $K$. Ако $X$ е допирната точка с $AB$ на външновписаната окръжност срещу върха $C$ за триъгълник $ABC$, то да се докаже, че правите $XC,KD$ и $IF$ се пресичат в една точка.

Решение. Първо ще докажем, че $CX$ и $IF$ се пресичат върху вписаната окръжност. Нека $CX$ пресича вписаната окръжност в точка $P$. Тогава, ако разгледаме хомотетия $h$ с център $C$, която изпраща вписаната окръжност във външновписаната окръжност срещу върха $C$, то $h(P)=X$ и следователно ако $t$ е допирателната през $P$ към вписаната окръжност, то имаме $h(t)=AB$. Последното означава, че $t\parallel AB$ и т.к $IP\perp t$ и $IF\perp AB$, то получаваме, че $P\in IF$.

От друга страна, ако $KD$ пресича вписаната окръжност за втори път в точка $Q$, то имаме $\Varangle EQD=180^\circ-\Varangle EFD=90^\circ+\frac{\Varangle ACB}{2}$. Имаме, че $\Varangle EKD=\frac{\Varangle ACB}{2}$, откъдето следва, че $\Varangle QEF=90^\circ$, т.е $Q=P$. Така получаваме, че правите $KD,CX$ и $IF$ се пресичат в точка $P$.

Оценяване. (6 точки) 3т. за $CX\cap IF$ лежи на вписаната окръжност и 3т. за $P\in KD$.

Задача 12.3. Да се реши в естествени числа уравнението:
$$
\sqrt[n]{m}+\sqrt[m]{n}=2+\frac{2}{mn(m+n)^{\frac1m+\frac1n}}.
$$

Решение. Уравнението няма решение в естествени числа. Ще докажем, че за всички двойки естествени числа $(m,n)$, за които имаме $\{\sqrt[n]{m}\}>0$, то е в сила неравенството $\{\sqrt[n]{m}\}\geq\frac1{mn}$. Нека $a=\lfloor\sqrt[n]{m}\rfloor$ и $x=\{\sqrt[n]{m}\}$. Тогава имаме, че
$$
m-a^n=x\sum_{i=0}^{n-1}m^{\frac in}a^{n-1-i}
$$
и следователно
$$
x\geq\frac1{nm^{\frac{n-1}{n}}}\geq\frac1{nm}.
$$
Така получаваме, че $\sqrt[n]{m}+\sqrt[m]{n}\geq2+\frac2{mn}$, откъдето задачата следва.

Оценяване. (7 точки) 2т. за формулиране на вярно неравенство за дробната част на $\sqrt[n]{m}$, 4т. за доказателството му и 1т. за довършване.

Задача 12.4. Ще наричаме едно множество $S$ от точки в равнината есенно, ако разстоянието между всеки две точки от $S$ е най-много 1. С $f(n,d)$ означаваме най-голямото цяло число, такова че за всяко есенно множество $S$ от $3n$ точки в равнината, съществува кръг с диаметър $d$, който съдържа поне $f(n,d)$ точки от $S$.

Да се докаже, че съществува $\varepsilon>0$ (независещо от $n$), за което за всички $d\in(1-\varepsilon,1)$ стойността на $f(n,d)$ не зависи от $d$ и да се определи тази стойност като функция на $n$.

Решение. Първо ще докажем, че съществува $d<1$, за което винаги можем да намерим окръжност с диаметър $d$ която съдържа $n$ точки от дадените. Нека $M$ е множество от $3n$ точки с диаметър 1 и нека за всяка точка $A\in M$ разгледаме кръг $D_A$ с център $A$ и радиус $\frac1{\sqrt3}$.

Лема. Ако $A,B,C\in M$, то $D_A\cap D_B\cap D_C\ne\emptyset$.

Доказателство. Ако точките $A,B,C$ образуват остроъгълен триъгълник, то някой от ъглите на триъгълника е с мярка между $60^\circ$ и $90^\circ$. Нека БОО това е $\Varangle ACB$. Тогава за радиуса $R$ на описаната около $ABC$ окръжност, получаваме от синусова теорема, че $2R=\frac{AB}{\sin\Varangle ACB}\leq\frac2{\sqrt3}$, откъдето следва, че $R\leq\frac1{\sqrt3}$ и следователно центърът на описаната окръжност $O\in D_A\cap D_B\cap D_C$.

Ако точките $A,B,C$ образуват тъпоъгълен или правоъгълен триъгълник, то нека БОО $\Varangle ACB\geq90^\circ$. Нека $M$ е средата на $AB$. Ясно е, че $MA\leq\frac12<\frac1{\sqrt3}$ и аналогично $MB<\frac1{\sqrt3}$, т.е $M\in D_A\cap D_B$. От друга страна
$$
MC^2=\frac{2AC^2+2CB^2-AB^2}{4}\leq\frac{AB^2}{4}\leq\frac14,
$$
където първото неравенство следва от това, че $\Varangle ACB\geq90^\circ$. Следователно $MC\leq\frac12<\frac1{\sqrt3}$ и следователно $M\in D_C$, т.е наистина $M\in D_A\cap D_B\cap D_C$, с което лемата е доказана.

От теорема на Хели, приложена за множеството от кръгове $\{D_A:A\in M\}$, следва, че т.к всеки три от тях се пресичат, то $\bigcap_{A\in M}D_A\ne\emptyset$. Нека $P\in\bigcap_{A\in M}D_A$. Тогава кръгът $D_P$ с център $P$ и радиус $\frac1{\sqrt3}$ съдържа $M$.

Сега ще покажем, че всеки кръг $K$ с радиус $\frac1{\sqrt3}$ може да бъде покрит от 3 кръга с радиус $\frac12$. Нека точките $A',B',C'\in D_P$ са такива, че $A'B'C'$ е равностранен триъгълник. Тогава $A'B'=B'C'=C'A'=1$ и кръговете с диаметри $A'B',A'C',B'C'$ покриват $D_P$.

Аналогично кръгове $k_1,k_2,k_3$ с центрове средите на $A'B',A'C',B'C'$ и диаметър $d$ "почти" покриват $K$. Наистина множеството $K\setminus(k_1\cup k_2\cup k_3)$ се състои от 3 еднакви фигури, които ще наричаме антилуни поради визуалната прилика. Да отбележим също, че можем да построим антилуни за всеки $,B,C$, лежащи на контура на $K$ и образуващи равностранен триъгълик.

Нека сега да допуснем противното, а именно, че за всяко $d<1$ можем да изберем множество от $3n$ точки , така че никои $n$ от тях да не лежат в кръг с диаметър $d$. Да забележим, че ако в трите антилуни съответсващи на някои $,B,C$ няма точки от , то поне един от $k_1,k_2,k_3$ съдържа $n$ точки. Оттук можем да заключим, че във всяка тройка съответни антилуни има точка от $M$.

Нека да фиксираме $N$ нечетно и да вземем $3N$ антилуни $A_1,\ldots,A_{3N}$ през равни ъгли. Нека също така сме избрали $d<1$, така че $\operatorname{dist}(A_1,A_{N+2})>1$. Това е възможно, защото когато $d$ клони към 1, диаметърът на антилуните сходи към 0. Тъй като разстоянието между две точки, образуващи ъгъл $2\pi(N+1)/3N$, е повече от 1, то от неравенството на триъгълника следва, че съществува $d<1$, за което $\operatorname{dist}(A_1,A_{N+2})>1$.

От предходните ни разсъждения следва, че за всяко $i$ поне една от антилуните $A_i,A_{N+i},A_{2N+i}$ съдържа точка от $M$. Нека БОО $A_1$ съдържа точка от . Следователно, антилуни те $A_{N+2},\ldots,A_{2N}$ не съдържат точки от . Нека $A_m$ е първата антилуна след $A_{2N}$, която съдържа точка. Следователно антилуните $A_{m-(2N-1)},\ldots,A_{m-1}$ не съдържат точки от $M$. От това следва, че антилуни $A_m,\ldots,A_{m-(2N+1)}$ съдържат точки. (Антилуната $A_{m-2N}$ е единствената, за която не знаем със сигурност дали съдържа точка). Нека $B_i$ е точката от $M$ в антилуна $A_{m-1+i}$ и нека $C_i$ е центърът на дъгата на $A_{m-1+i}$

Сега можем да изберем $N$ достатъчно голямо и $d=d'$ достатъчно близко до 1 така, че:
$$
\begin{gathered}
\operatorname{dist}(C_1,B_1)\leq\operatorname{diam}(A_m)<0.01\\
\operatorname{dist}((C_1+C_{N+1})/2,B_{\lfloor(N+1)/2\rfloor})\leq\operatorname{diam}(A_m)+2\pi/3N<0.01\\
\operatorname{dist}(C_{N+1},B_N)\leq\operatorname{diam}(A_m)+2\pi/3N<0,01
\end{gathered}
$$

Следователно, множеството $M$ лежи изцяло в сечението $P$ на и трите кръга $k_1,k_2,k_3$ с центрове $C_1,C_{(N+1)/2},C_N$ и диаметър 1.01. Но сега лесно се забелязва, че $P$ може да се впише в окръжност $K'$ с радиус $d''/\sqrt3<1/\sqrt3$. Но тогава ${}^{\prime}$ се покрива (съответно и ) от 3 кръга с диаметър $d''$. Следователно, за $d=\max(d',d'')$ със сигурност съществува кръг покриващ поне $n$ точки от $M$.

Остана да покажем, че за всяко $d<1$ и всяко $n\in\mathbb N$ съществува множество $M$ от $3n$ точки, никои $n+1$ от които не принадлежат на един кръг с диаметър $d$. Построяваме $n$ еднакво ориентирани равностранни триъгълника със страна $(1+d)/2$, така че разстоянието между всеки 2 съответни върха е по-малко от $(1-d)/10$. Лесно се проверя, че тази конструкция изпълнява необходимите условия.

Оценяване. (7 точки) 2т. за доказателство, че $S$ се покрива от 3 кръга с диаметър 1, 4т. за показване, че $S$ се покрива от 3 кръга с диаметър $d'<1$, 1т. за конструкция на множество от $3n$ точки, от което не може да бъдат покрити $n+1$ точки с кръг с диаметър $d$ за фиксирано $d<1$.
